% sections/en/5-experiments.tex -- structured skeleton per spec 6C (5.1-5.7 + instance-factor effects); numbers are placeholders.
\section{Computational experiments}\label{sec:experiments}

\draftnote{Every number in this section is a placeholder to be filled by \texttt{paper\_assets/scripts/make\_tables.py} from
\texttt{result/*.csv} under the 6B logging contract; every figure is a placeholder box to be replaced by opt-paper-scifig. All
directional statements are expectations until the claims are judged.}

\subsection{Experimental settings}\label{sec:settings}
The evaluation set is a fixed design set generated from an instance parameter table with the mean multiplicity $\bar N$ as scale
factor and three structural factors, multiplicity dispersion $V_N$, flexibility $F$ and processing-time heterogeneity $V_p$,
with $R = 8$ job types and $m = 10$ machines (Table~\ref{tab:instances}). The main tier is a $3\times3\times3$ full
factorial at $\bar N = 10$; the small tier carries exact references from CP-SAT \cite{ortools} at 1800 s; the large tier
extrapolates the multiplicity beyond the training range; the Brandimarte instances \cite{brandimarte1993routing} are the
$N_r = 1$ special case. Instances are stored as files and released with the paper; reproduction relies on them and on
\NumRuns{} independent runs, not on random seeds. Table~\ref{tab:hyper} lists the hyperparameters and
Table~\ref{tab:baselines} the baselines, one strong representative per family, each at the recommended configuration and the
same wall-clock budget. Metrics are the makespan, the optimality gap against the exact or reference solution, the fluid gap
$(C_{\max} - L^F)/L^F$ for cross-tier comparison, the feasibility rate, the area under the quality--time curve at 1, 10 and 60 s
checkpoints, and the number of exact decodings and the time to target. Significance is assessed by paired Wilcoxon tests on
per-instance run means ($p < 0.05$) and, with four or more methods, by the Friedman--Nemenyi critical-difference diagram.

\begin{table*}
\centering\small
\caption{Evaluation instance set.}\label{tab:instances}
% generated by paper_assets/scripts/make_tables.py -- do not edit by hand

\begin{tabularx}{\linewidth}{@{}l>{\hsize=1.000\hsize\raggedright\arraybackslash}X>{\hsize=0.900\hsize\raggedright\arraybackslash}Xl>{\hsize=1.100\hsize\raggedright\arraybackslash}X@{}}
\toprule
\textbf{\bi{档位}{Tier}} & \textbf{\bi{规模因子 $\bar N$}{Scale factor $\bar N$}} & \textbf{\bi{结构因子}{Structural factors}} & \textbf{\bi{算例数}{Instances}} & \textbf{\bi{用途}{Purpose}} \\
\midrule
Small & 2, 4 & \bi{各取低 / 高两端}{low / high ends} & 16 & \bi{精确参考解校准真实 gap}{exact reference solutions calibrate the true gap} \\
Main & 10 & \bi{3 $\times$ 3 $\times$ 3 全因子}{$3\times3\times3$ full factorial} & 27 & \bi{主对比 + 因子效应分析}{main comparison and factor-effect analysis} \\
Large & 25, 50 & \bi{各取低 / 高两端}{low / high ends} & 16 & \bi{多重度外推（商图表示直接求解）}{multiplicity extrapolation (quotient representation solved directly)} \\
Benchmark & \bi{Brandimarte mk01–mk10（$N_r = 1$ 退化为标准 FJSP）}{Brandimarte mk01--mk10 ($N_r = 1$, standard FJSP)} & \bi{外部给定}{given} & 10 & \bi{公开基准（补充，不替代设计集）}{public benchmark (supplementary to the design set)} \\
\bottomrule
\end{tabularx}

\end{table*}

\begin{table}
\centering\small
\caption{Hyperparameters (draft values; final values from the confirmed training-parameter table).}\label{tab:hyper}
% generated by paper_assets/scripts/make_tables.py -- do not edit by hand

\begin{tabularx}{\linewidth}{@{}l>{\hsize=0.850\hsize\raggedright\arraybackslash}X>{\hsize=1.150\hsize\raggedright\arraybackslash}X@{}}
\toprule
\textbf{\bi{类别}{Group}} & \textbf{\bi{参数}{Parameter}} & \textbf{\bi{取值}{Value}} \\
\midrule
\bi{训练}{Training} & \texttt{train.paradigm} & \bi{generalist（排序器离线训练、冻结后零样本）}{generalist (ranker trained offline, frozen, zero-shot)} \\
\bi{训练}{Training} & \texttt{instances\_\allowbreak{}per\_\allowbreak{}epoch} & \bi{256 个随机算例，FRES-lite 自举 32\,768 个候选-标签对}{256 random instances; 32\,768 candidate--label pairs bootstrapped by FRES-lite} \\
\bi{训练}{Training} & \texttt{epochs} /\allowbreak{} \texttt{plateau\_\allowbreak{}patience} & 60 / 10 \\
\bi{训练}{Training} & \texttt{wall\_\allowbreak{}budget\_\allowbreak{}s} /\allowbreak{} \texttt{eval\_\allowbreak{}every} /\allowbreak{} \texttt{env\_\allowbreak{}batch} & \bi{待实测}{to be measured} \\
\bi{算法}{Optimizer} & \texttt{lr} & \bi{1e-3，余弦退火}{1e-3, cosine annealing} \\
\bi{算法}{Optimizer} & \texttt{batch\_\allowbreak{}size} /\allowbreak{} \texttt{loss} & \bi{512 / pairwise ranking（logistic）}{512 / pairwise ranking (logistic)} \\
\bi{网络}{Network} & \texttt{embed\_\allowbreak{}dim} /\allowbreak{} \texttt{n\_\allowbreak{}layers} & \bi{128 / 3（异构消息传递）}{128 / 3 (heterogeneous message passing)} \\
\bi{宿主}{Host EA} & \texttt{population} /\allowbreak{} \texttt{offspring} & 100 / 100 \\
\bi{宿主}{Host EA} & $B$ / $K_L$ / $K_H$ / $K_R$ & 96 / 10 / 5 / 3 \\
\bi{宿主}{Host EA} & \bi{配额迁移步长 $b$ / 优先键扰动幅度}{quota step $b$ / priority-key perturbation} & 1 / 0.1 \\
\bottomrule
\end{tabularx}

\end{table}

\begin{table*}
\centering\small
\caption{Baselines and their configuration.}\label{tab:baselines}
% generated by paper_assets/scripts/make_tables.py -- do not edit by hand

\begin{tabularx}{\linewidth}{@{}>{\hsize=0.700\hsize\raggedright\arraybackslash}Xl>{\hsize=1.300\hsize\raggedright\arraybackslash}X@{}}
\toprule
\textbf{\bi{方法}{Method}} & \textbf{\bi{类别}{Class}} & \textbf{\bi{实现与配置}{Implementation and configuration}} \\
\midrule
\bi{HGA（job-indexed）}{HGA (job-indexed)} & \bi{元启发式}{metaheuristic} & \bi{文献推荐参数（双层编码遗传算法）；种群 100、POX 交叉 0.8、变异 0.15、精英 10\%；同墙钟；每算例 3 次重启}{recommended settings of the two-layer encoded GA; population 100, POX crossover 0.8, mutation 0.15, elitism 10\%; same wall clock; 3 restarts per instance} \\
\bi{关键块禁忌搜索（TS）}{Critical-block tabu search (TS)} & \bi{元启发式}{metaheuristic} & \bi{Mastrolilli–Gambardella 型，自实现并在 mk 基准核对；禁忌长度 12、关键块邻域 + 关键工序改派；同墙钟；3 次重启}{Mastrolilli--Gambardella type, reimplemented and checked on the mk benchmark; tenure 12, critical-block neighbourhood with reassignment; same wall clock; 3 restarts} \\
\bi{流体进化算法（FMAE-s）}{Fluid-assisted EA (FMAE-s)} & \bi{元启发式}{metaheuristic} & \bi{流体主从进化算法的无换模静态版（复现）；原文推荐参数；流体比例只作初始化与规则；同墙钟；3 次重启}{static, setup-free version of the fluid master--slave EA (reimplementation); published settings; fluid ratios used only for initialization and rules; same wall clock; 3 restarts} \\
DSLFMAE & \bi{学习型}{learning-based} & \bi{PPO 辅助流体进化算法（复现或机制等价复现，标注 reimplementation）；原文推荐配置；与本方法同等训练预算，3 run}{PPO-assisted fluid EA (reimplementation, marked as such); published configuration; same training budget as ours, 3 runs} \\
HGNN-DRL & \bi{学习型}{learning-based} & \bi{官方代码（异构图神经网络 + PPO 构造式调度）；推荐配置；贪心 + 采样两口径；同等训练预算，3 run}{official code (heterogeneous GNN with PPO, constructive); recommended configuration; greedy and sampling decoding; same training budget, 3 runs} \\
CP-SAT & \bi{参考解}{reference} & \bi{OR-Tools；1800 s（small）/ 与元启发式同墙钟（main，anytime 参考）}{OR-Tools; 1800 s (small) / same wall clock as the metaheuristics (main, anytime reference)} \\
\bottomrule
\end{tabularx}

\end{table*}

\subsection{Comparison with baselines}\label{sec:comparison}
This subsection asks whether FRES is better than the metaheuristic and learning-based baselines, and whether the difference is
significant. Table~\ref{tab:main} reports the small and main tiers at the 10 s and 60 s checkpoints; Fig.~\ref{fig:main}
shows the gap distributions and the critical-difference diagram. \pending{FRES-lite alone reaches a gap no higher than HGA
and FMAE-s, and the full method improves on the strongest baseline by \ImpvOverBest{} on the main tier.}

\begin{table*}
\centering\small
\caption{Main results: mean $\pm$ std over \NumRuns{} runs; best in bold; $\dagger$ marks significance over the runner-up (Wilcoxon, $p < 0.05$). S = small tier (gap to the exact solution), M = main tier (gap to the reference).}\label{tab:main}
% generated by paper_assets/scripts/make_tables.py -- do not edit by hand

\begin{tabularx}{\linewidth}{@{}>{\hsize=1.500\hsize\raggedright\arraybackslash}X>{\hsize=0.800\hsize\raggedright\arraybackslash}X>{\hsize=0.800\hsize\raggedright\arraybackslash}X>{\hsize=1.000\hsize\raggedright\arraybackslash}X>{\hsize=1.000\hsize\raggedright\arraybackslash}X>{\hsize=0.800\hsize\raggedright\arraybackslash}X>{\hsize=1.100\hsize\raggedright\arraybackslash}X@{}}
\toprule
\textbf{\bi{方法}{Method}} & \textbf{\bi{S-gap（\%）}{S gap (\%)}} & \textbf{\bi{S-可行率}{S feas.}} & \textbf{\bi{M-gap 10 s（\%）}{M gap 10 s (\%)}} & \textbf{\bi{M-gap 60 s（\%）}{M gap 60 s (\%)}} & \textbf{\bi{M-可行率}{M feas.}} & \textbf{\bi{解码次数（$\times10^{3}$）}{Decodings ($\times10^{3}$)}} \\
\midrule
\bi{HGA（job-indexed）}{HGA (job-indexed)} & \phpm & \ph & \phpm & \phpm & \ph & \ph \\
\bi{关键块禁忌搜索（TS）}{Critical-block tabu search (TS)} & \phpm & \ph & \phpm & \phpm & \ph & \ph \\
\bi{流体进化算法（FMAE-s）}{Fluid-assisted EA (FMAE-s)} & \phpm & \ph & \phpm & \phpm & \ph & \ph \\
\bi{DSLFMAE}{DSLFMAE} & \phpm & \ph & \phpm & \phpm & \ph & \ph \\
\bi{HGNN-DRL}{HGNN-DRL} & \phpm & \ph & \phpm & \phpm & \ph & \ph \\
\bi{CP-SAT（同墙钟）}{CP-SAT (same wall clock)} & \phpm & \ph & \phpm & \phpm & \ph & \ph \\
\bi{\textbf{FRES}}{\textbf{FRES}} & \phb & \phb & \phb & \phb & \phb & \phb \\
\bottomrule
\end{tabularx}

\end{table*}

\begin{figure}
\figplace{gap box plots per method on the main tier at 60 s, one panel per tier, with Wilcoxon significance brackets; inset critical-difference diagram}{\texttt{result/*/eval\_results.csv} (gap, method, tier); $p$ values and ranks from \texttt{result/stats\_summary.csv}}{opt-paper-scifig route A (matplotlib), column width 88\,mm}{52mm}
\caption{Gap distributions and critical-difference diagram on the main tier.}\label{fig:main}
\end{figure}

\subsection{Ablation and ranking study}\label{sec:ablation}
This subsection asks whether the gain comes from the three mechanisms and whether the ranker reduces wasted decodings.
Variants A1--A5 remove, one at a time, the fluid quota, the quota compression, the ranker, the residual scoping and the
safeguard (Table~\ref{tab:ablation}). Fig.~\ref{fig:ablation} shows the convergence of the variants, Fig.~\ref{fig:recall}
the top-10 recall of the ranker against the fluid heuristic and the generic ranker at equal candidate budget, and
Fig.~\ref{fig:safeguard} the share of accepted offspring that came from each source of the union. \pending{The full method
decodes \OursDecodeReduction{} fewer candidates than FRES-lite at no loss of quality.}

\begin{table}
\centering\small
\caption{Ablation and mechanism metrics on the main tier.}\label{tab:ablation}
% generated by paper_assets/scripts/make_tables.py -- do not edit by hand

\begin{tabularx}{\linewidth}{@{}>{\hsize=1.000\hsize\raggedright\arraybackslash}Xllll@{}}
\toprule
\textbf{\bi{变体}{Variant}} & \textbf{\bi{gap（\%）}{gap (\%)}} & \textbf{\bi{可行率}{feas.}} & \textbf{top-10 recall} & \textbf{\bi{解码次数（$\times10^{3}$）}{Decodings ($\times10^{3}$)}} \\
\midrule
\bi{完整方法（FRES）}{Full method (FRES)} & \phb & \phb & \phb & \phb \\
\bi{A1 无流体配额}{A1 no fluid quota} & \phpm & \ph & \ph & \ph \\
\bi{A2 无配额压缩}{A2 no quota compression} & \phpm & \ph & \ph & \ph \\
\bi{A3 无排序器（FRES-lite）}{A3 no ranker (FRES-lite)} & \phpm & \ph & \ph & \ph \\
\bi{A4 通用排序器}{A4 generic ranker} & \phpm & \ph & \ph & \ph \\
\bi{A5 无保底}{A5 no safeguard} & \phpm & \ph & \ph & \ph \\
\bottomrule
\end{tabularx}

\end{table}

\begin{figure}
\figplace{best-so-far gap versus wall clock for FRES and A1--A5, mean line with $\pm$ std band, log time axis}{\texttt{result/*/curve.csv} aggregated over runs}{opt-paper-scifig route A, column width 88\,mm}{50mm}
\caption{Convergence of the ablation variants.}\label{fig:ablation}
\end{figure}

\begin{figure}
\figplace{top-10 recall of the true improving action versus the number of decoded candidates, three curves (residual ranker, fluid heuristic, generic ranker), same candidate pool}{\texttt{result/*/log.csv} ranking metrics and \texttt{operator\_trace.csv}}{opt-paper-scifig route A, column width 88\,mm}{50mm}
\caption{Ranking quality at equal candidate budget.}\label{fig:recall}
\end{figure}

\begin{figure}
\figplace{stacked bars of the source (learned / fluid heuristic / random) of the accepted offspring over generations, one panel per $V_N$ level}{\texttt{operator\_trace.csv} (best\_from, generation) joined with \texttt{index.csv}}{opt-paper-scifig route A, column width 88\,mm}{48mm}
\caption{Source of the accepted offspring under the safeguarded union.}\label{fig:safeguard}
\end{figure}

\subsection{Parameter sensitivity}\label{sec:sensitivity}
This subsection asks whether the preselection parameters need tuning. Fig.~\ref{fig:sens} scans the pool size $B$ and the
learner share $K_L$ on the validation tier. \pending{Performance is flat over the scanned grid.}

\begin{figure}
\figplace{heat map of the main-tier gap over $B \in \{48, 96, 192\}$ and $K_L \in \{8, 16, 24\}$, cell annotated with the mean, chosen setting outlined}{\texttt{result/sensitivity.csv}}{opt-paper-scifig route A, column width 88\,mm}{48mm}
\caption{Sensitivity to the candidate-pool size and the learner share.}\label{fig:sens}
\end{figure}

\subsection{Instance-factor effects}\label{sec:factor}
This subsection asks on which problem characteristics the method is stronger or weaker, and whether the theoretical bound is
supported. It is distinct from the parameter sensitivity above: here the factors belong to the problem, not to the algorithm.
Fig.~\ref{fig:rounding} compares the empirical rounding gap $G_Q/L^F$ with the bound $\bar B/L^F$ of
Proposition~\ref{prop:rounding} across $\bar N$ and $V_p$ (Conjecture~\ref{conj:rounding}); Table~\ref{tab:factor} and
Fig.~\ref{fig:factor} give the main effects and two-factor interactions of the gain over the strongest baseline
(Conjecture~\ref{conj:gain}). Error bars are between-run standard deviations, not within-cell replicates.

\begin{table*}
\centering\small
\caption{Instance-factor main effects on the main tier (27 instances with distinct parameters, 9 per level; error is the between-run std over \NumRuns{} runs).}\label{tab:factor}
% generated by paper_assets/scripts/make_tables.py -- do not edit by hand

\begin{tabularx}{\linewidth}{@{}>{\hsize=0.900\hsize\raggedright\arraybackslash}Xll>{\hsize=1.100\hsize\raggedright\arraybackslash}Xl@{}}
\toprule
\textbf{\bi{因子}{Factor}} & \textbf{\bi{水平}{Level}} & \textbf{\bi{gap（\%）}{gap (\%)}} & \textbf{\bi{相对最强基线的增益（pp）}{Gain over the strongest baseline (pp)}} & \textbf{\bi{解码次数（$\times10^{3}$）}{Decodings ($\times10^{3}$)}} \\
\midrule
\bi{多重度离散度 $V_N$}{Multiplicity dispersion $V_N$} & 0.0 & \phpm & \ph & \ph \\
\bi{多重度离散度 $V_N$}{Multiplicity dispersion $V_N$} & 0.5 & \phpm & \ph & \ph \\
\bi{多重度离散度 $V_N$}{Multiplicity dispersion $V_N$} & 1.0 & \phpm & \ph & \ph \\
\bi{柔性度 $F$}{Flexibility $F$} & 0.3 & \phpm & \ph & \ph \\
\bi{柔性度 $F$}{Flexibility $F$} & 0.55 & \phpm & \ph & \ph \\
\bi{柔性度 $F$}{Flexibility $F$} & 0.8 & \phpm & \ph & \ph \\
\bi{工时异质度 $V_p$}{Processing-time heterogeneity $V_p$} & 0.2 & \phpm & \ph & \ph \\
\bi{工时异质度 $V_p$}{Processing-time heterogeneity $V_p$} & 0.5 & \phpm & \ph & \ph \\
\bi{工时异质度 $V_p$}{Processing-time heterogeneity $V_p$} & 0.8 & \phpm & \ph & \ph \\
\bottomrule
\end{tabularx}

\end{table*}

\begin{figure}
\figplace{two panels: empirical $G_Q/L^F$ and the bound $\bar B/L^F$ versus $\bar N$ (log axis, small / main / large tiers) and versus $V_p$ (main tier), markers with between-run error bars}{\texttt{relaxation\_diag.csv} (rounding\_gap, rounding\_bound, fluid\_lb) joined with \texttt{index.csv}}{opt-paper-scifig route A, column width 88\,mm}{48mm}
\caption{Empirical rounding gap against its theoretical bound.}\label{fig:rounding}
\end{figure}

\begin{figure*}
\figplace{left: main-effect plot of the gain over the strongest baseline for $V_N$, $F$, $V_p$ with error bars; right: two-factor interaction heat maps}{\texttt{eval\_results.csv} joined with \texttt{index.csv} on instance\_id}{opt-paper-scifig route A, two-column width 178\,mm}{52mm}
\caption{Factor effects of the gain over the strongest baseline.}\label{fig:factor}
\end{figure*}

\subsection{Training and search behaviour}\label{sec:training}
This subsection asks whether the search and the ranker training are stable. Fig.~\ref{fig:anytime} shows the best-so-far gap
on the fixed validation set against wall clock for all methods and the ranker's training curve.

\begin{figure}
\figplace{left: anytime curve of the validation-set gap versus wall clock for all methods, mean $\pm$ std; right: ranking loss and top-10 recall per epoch for the ranker over \NumRuns{} training runs}{\texttt{curve.csv} and \texttt{log.csv}}{opt-paper-scifig route A, column width 88\,mm}{50mm}
\caption{Anytime behaviour and ranker training curves.}\label{fig:anytime}
\end{figure}

\subsection{Generalization to higher multiplicity}\label{sec:generalization}
This subsection asks whether the method works only within the training multiplicity range. Table~\ref{tab:generalization}
solves the large tier ($\bar N \in \{25, 50\}$) directly with the frozen ranker; Fig.~\ref{fig:gen} plots the gap against
$\bar N$ with the training range shaded. \pending{The gap of FRES is flat in $\bar N$ while that of the job-indexed hosts
grows.}

\begin{table}
\centering\small
\caption{Multiplicity extrapolation beyond the training range ($\bar N \le 12$).}\label{tab:generalization}
% generated by paper_assets/scripts/make_tables.py -- do not edit by hand

\begin{tabularx}{\linewidth}{@{}>{\hsize=1.100\hsize\raggedright\arraybackslash}X>{\hsize=1.000\hsize\raggedright\arraybackslash}X>{\hsize=0.900\hsize\raggedright\arraybackslash}X>{\hsize=1.000\hsize\raggedright\arraybackslash}X>{\hsize=0.900\hsize\raggedright\arraybackslash}X@{}}
\toprule
\textbf{\bi{方法}{Method}} & \textbf{\bi{$\bar N = 25$ gap（\%）}{$\bar N = 25$ gap (\%)}} & \textbf{\bi{$\bar N = 25$ 可行率}{$\bar N = 25$ feas.}} & \textbf{\bi{$\bar N = 50$ gap（\%）}{$\bar N = 50$ gap (\%)}} & \textbf{\bi{$\bar N = 50$ 可行率}{$\bar N = 50$ feas.}} \\
\midrule
\bi{HGA（job-indexed）}{HGA (job-indexed)} & \phpm & \ph & \phpm & \ph \\
\bi{流体进化算法（FMAE-s）}{Fluid-assisted EA (FMAE-s)} & \phpm & \ph & \phpm & \ph \\
\bi{A2 无配额压缩}{A2 no quota compression} & \phpm & \ph & \phpm & \ph \\
\bi{\textbf{FRES}}{\textbf{FRES}} & \phb & \phb & \phb & \phb \\
\bottomrule
\end{tabularx}

\end{table}

\begin{figure}
\figplace{gap versus $\bar N$ for FRES, HGA, FMAE-s and A2, training range shaded, log $x$ axis}{\texttt{eval\_results.csv} (tier, gap) joined with \texttt{index.csv} (scale)}{opt-paper-scifig route A, column width 88\,mm}{48mm}
\caption{Gap against the mean multiplicity.}\label{fig:gen}
\end{figure}

\subsection{Case study}\label{sec:case}
This subsection asks whether the schedules are reasonable and whether the quotas stay close to the fluid flow.
Fig.~\ref{fig:case} shows a Gantt chart of one main-tier instance next to the integer quota and the fluid ratio per
type-operation.

\begin{figure*}
\figplace{left: Gantt chart of the FRES schedule on one main-tier instance, copies of a type in one hue, critical path outlined; right: per type-operation bars of $n_{rki}/N_r$ against $q^*_{rki}/N_r$}{schedule export of one instance; \texttt{relaxation\_diag.csv}}{opt-paper-scifig route A, two-column width 178\,mm}{60mm}
\caption{Schedule and quota--fluid comparison on one instance.}\label{fig:case}
\end{figure*}
